% =====================================================================
%  项目一图流海报 —— abb-offline-coder  (A4 横版，答辩展板/路演用)
%  比赛赛项：人工智能创新赛   队伍：人工智能创新001
%  编译：xelatex poster.tex
% =====================================================================
\documentclass[12pt]{ctexart}
\usepackage[paperwidth=297mm,paperheight=210mm,margin=9mm]{geometry}
\setCJKmainfont{Noto Serif CJK SC}
\setCJKsansfont{Noto Sans CJK SC}
\setCJKmonofont{Noto Sans Mono CJK SC}
\usepackage{tikz}
\usetikzlibrary{positioning,arrows.meta,calc,fit,backgrounds}
\usepackage{xcolor}
\usepackage{amssymb}
\usepackage{pifont}
\pagestyle{empty}
\setlength{\parindent}{0pt}

\definecolor{abbred}{RGB}{217,74,31}
\definecolor{abbdark}{RGB}{24,24,24}
\definecolor{abbgray}{RGB}{95,95,95}
\definecolor{abblight}{RGB}{246,242,240}
\definecolor{boxblue}{RGB}{226,236,246}
\definecolor{okgreen}{RGB}{34,140,64}
\definecolor{linegray}{RGB}{205,205,205}

\newcommand{\hd}[1]{{\sffamily\bfseries\color{abbred}#1}}

\begin{document}
\begin{tikzpicture}[x=1cm,y=1cm,
  panel/.style={draw=abbred,line width=0.9pt,rounded corners=4pt,fill=white,
    inner sep=7pt,align=left,font=\footnotesize},
  pstep/.style={draw=abbdark,line width=0.7pt,rounded corners=3pt,fill=white,
    align=center,font=\scriptsize\sffamily,minimum height=1.05cm,inner sep=2pt},
  pio/.style={draw=abbred,line width=1pt,fill=abblight,rounded corners=3pt,
    align=center,font=\scriptsize\sffamily\bfseries,minimum height=1.05cm,inner sep=2pt},
  lyr/.style={draw=abbdark,line width=0.6pt,rounded corners=2pt,fill=abblight,
    align=center,font=\scriptsize\sffamily,minimum height=0.62cm,inner sep=2pt},
  arr/.style={-{Latex[length=2mm]},line width=0.8pt,abbred},
]
\useasboundingbox (0,0) rectangle (27.9,19.2);

% ---------------- 顶部标题带 ----------------
\fill[abbred] (0,17.2) rectangle (27.9,19.2);
\fill[abbdark] (0,17.0) rectangle (27.9,17.2);
\node[anchor=west,font=\sffamily\bfseries,text=white] at (0.5,18.5)
  {\fontsize{30}{32}\selectfont abb-offline-coder};
\node[anchor=west,font=\sffamily,text=white] at (0.6,17.65)
  {\large 面向 ABB 喷涂机器人的离线 AI 编程助手\, ·\, “把中文一句话变成 RAPID 喷涂程序”};
\node[anchor=east,font=\sffamily\bfseries,text=white,align=right] at (27.5,18.45)
  {人工智能创新赛};
\node[anchor=east,font=\sffamily,text=white,align=right] at (27.5,17.7)
  {队伍：人工智能创新001\quad|\quad MIT 开源};

% ---------------- 流水线带（核心）----------------
\node[anchor=north west,font=\sffamily\bfseries\color{abbred}] at (0.1,16.85)
  {核心流水线：中文需求 → 可导入 RobotStudio 的 .mod 程序（全离线 / CPU）};
\def\py{15.0}        % 流水线行 y 中心
\node[pio,text width=2.7cm]   at (1.9,\py)  (q)  {中文需求\\\tiny 自然语言};
\node[pstep,text width=2.7cm] at (5.5,\py)  (rw) {\ding{172}~查询改写\\\tiny 中→英术语+分类};
\node[pstep,text width=2.7cm] at (9.1,\py)  (re) {\ding{173}~混合检索\\\tiny 向量+BM25};
\node[pstep,text width=2.7cm] at (12.7,\py) (cb) {\ding{174}~上下文装配\\\tiny 手册片段注入};
\node[pstep,text width=2.7cm] at (16.3,\py) (lm) {\ding{175}~本地 LLM\\\tiny Qwen2.5-Coder};
\node[pstep,text width=2.7cm] at (19.9,\py) (pp) {\ding{176}~校验后处理\\\tiny 11类规则+错误码};
\node[pio,text width=2.7cm]   at (23.5,\py) (o)  {\ding{177}~.mod 输出\\\tiny +Pack\&Go 交付};
\foreach \a/\b in {q/rw,rw/re,re/cb,cb/lm,lm/pp,pp/o} \draw[arr] (\a) -- (\b);
\node[anchor=west,font=\scriptsize\color{abbgray}] at (25.0,\py) {端到端\\20–35s};
\draw[linegray,line width=0.5pt] (0,13.85) -- (27.9,13.85);

% ================= 三栏 =================
\def\cyA{13.35}  % 各栏顶部 y
% ---------- 左栏 ----------
\node[panel,text width=8.0cm,anchor=north west] at (0.1,\cyA) (p1){%
  \hd{行业痛点：三高一隔离}\\[2pt]
  $\bullet$ RAPID 专有语言，喷涂还叠加笔刷/时序/标定/IO 配置——\textbf{门槛高、试错贵、依赖专家}\\
  $\bullet$ 喷涂车间多为\textbf{物理隔离、完全无网}——云端 AI 编程工具（Copilot/ChatGPT）\textbf{无法落地}，现场数据\textbf{不能外传}};
\node[panel,text width=8.0cm,anchor=north west,fill=abblight] at (0.1,\cyA-2.55) (p2){%
  \hd{解决方案与定位}\\[2pt]
  本地量化大模型 + 官方手册 RAG，在工业 PC（CPU、无 GPU）上把\textbf{一句中文}变成
  符合 ABB 规范、可导入 RobotStudio 的 \textbf{.mod} 程序；\textbf{数据全程不出厂}。};
\node[panel,text width=8.0cm,anchor=north west] at (0.1,\cyA-5.0) (p3){%
  \hd{七大创新点}\\[2pt]
  \textbf{1.} 完全离线工业代码生成\quad\textbf{2.} 官方手册 RAG 抑幻觉\\
  \textbf{3.} IRC5/IRC5P 双控制器自适应\quad\textbf{4.} 生成—校验闭环\\
  \textbf{5.} Pack\&Go 一键交付\quad\textbf{6.} 一键离线迁移\quad\textbf{7.} 安全护栏内建};

% ---------- 中栏 ----------
\node[anchor=north,font=\footnotesize\sffamily\bfseries\color{abbred}] at (13.6,\cyA-0.2) (aTitle){系统四层架构};
% 架构小图
\begin{scope}
\node[lyr,text width=7.2cm] at (13.6,\cyA-0.95) (aCli){CLI 层\,·\, gen/chat/kb/init/doctor};
\node[lyr,text width=7.2cm,fill=boxblue] at (13.6,\cyA-1.72) (aOrc){编排层 Agent\,·\,流水线+多轮对话};
\node[lyr,text width=2.25cm,minimum height=0.95cm,draw=abbred] at (11.45,\cyA-2.85) (aR){rag/\\\tiny 检索};
\node[lyr,text width=2.25cm,minimum height=0.95cm,draw=abbred] at (13.6,\cyA-2.85) (aL){llm/\\\tiny 推理};
\node[lyr,text width=2.25cm,minimum height=0.95cm,draw=abbred] at (15.75,\cyA-2.85) (aP){rapid/\\\tiny 校验};
\node[lyr,text width=7.2cm,fill=abblight] at (13.6,\cyA-3.75) (aKb){knowledge/\,·\,PDF/MOD/HTML→切片→ChromaDB};
\draw[arr] (aCli)--(aOrc);
\draw[arr] (aOrc) -- (aL); \draw[arr](aOrc)--(aR); \draw[arr](aOrc)--(aP);
\draw[arr] (aR) -- (aKb);
\end{scope}
\begin{scope}[on background layer]
\node[panel,fit=(aTitle)(aCli)(aOrc)(aR)(aL)(aP)(aKb),inner sep=8pt]{};
\end{scope}
\node[panel,text width=8.0cm,anchor=north west] at (9.6,\cyA-4.65) (m2){%
  \hd{IRC5\, vs\, IRC5P（同一句话，两种工艺）}\\[3pt]
  \textbf{IRC5}（通用）：\texttt{MoveL + SetDO doSprayOn} 手动开关枪\\
  \textbf{IRC5P}（Paint）：\texttt{PaintL/PaintC + brushdata} 原生工艺自动管时序\\
  控制器作为贯穿提示词/生成/校验/交付的\textbf{一等参数}，一键切换。};
\node[panel,text width=8.0cm,anchor=north west,fill=abblight] at (9.6,\cyA-7.55) (m3){%
  \hd{Pack\&Go 一键交付}\\[2pt]
  \texttt{.mod + BASE.sys + T\_ROB1.pgf + README} → 可直接拖入 RobotStudio
  加载的完整目录；自动处理控制器 ISO-8859-1/CRLF 编码细节。};

% ---------- 右栏 ----------
\node[panel,text width=8.05cm,anchor=north west] at (19.05,\cyA) (r1){%
  \hd{核心规格（实测 / 记录）}\\[3pt]
  \begin{tabular}{@{}ll@{}}
  单元测试 & \textbf{156} 通过 / 0.94s（实测）\\
  核心覆盖率 & 92\%–\textbf{100}\%（实测）\\
  代码规模 & $\approx$4506 行 / 37 文件\\
  知识库 & \textbf{7497} 片段 / 10 本手册\\
  嵌入模型 & bge-small-zh, 512维, 92MB\\
  本地 LLM & Qwen2.5-Coder-7B Q4 (4.5GB)\\
  端到端 & 20–35s（i5/16G，无 GPU）\\
  离线包 & $\approx$192MB（单文件迁移）\\
  \end{tabular}};
\node[panel,text width=8.05cm,anchor=north west] at (19.05,\cyA-5.35) (r2){%
  \hd{生成—校验闭环（实跑实录）}\\[2pt]
  对缺陷代码精确拦截，错误码可追溯：\\
  {\ttfamily\scriptsize\color{abbdark}[ERROR][PNT001] 缺 brushdata 形参\\
  $[$ERROR$][$IO001$]$ 信号不在 EIO.cfg 白名单}\\
  把现场才暴露的 40212 错误\textbf{左移到生成时}。};
\node[panel,text width=8.05cm,anchor=north west,fill=boxblue] at (19.05,\cyA-8.15) (r3){%
  \hd{开源 \& 在线演示}\\[2pt]
  \footnotesize 源码\,github.com/Hollis36/abb-offline-coder\\
  演示\,hollis36.github.io/abb-offline-coder\\
  \scriptsize Python 87\%\, ·\, MIT\, ·\, rag/ollama/qwen/chromadb/abb-robotics};

% ---------------- 底部带 ----------------
\fill[abblight] (0,0) rectangle (27.9,1.05);
\draw[abbred,line width=1.2pt] (0,1.05) -- (27.9,1.05);
\node[anchor=west,font=\scriptsize\sffamily\color{abbgray}] at (0.3,0.52)
  {人工智能创新赛 · 项目一图流 \quad|\quad 队伍：人工智能创新001 \quad|\quad abb-offline-coder v0.2};
\node[anchor=east,font=\scriptsize\sffamily\color{abbdark}] at (27.6,0.52)
  {完全离线 · 中文需求 → ABB RAPID · 本地 LLM + 官方手册 RAG · 工业 PC 断网可用};
\end{tikzpicture}
\end{document}
